\par\medskip
\noindent
11. \underline{Componentes conexas del modelo de Néron: relación de
dualidad y comportamiento asintótico}

\smallskip
\noindent
11.0. El objeto de la presente sección es estudiar el $k$-grupo finito
étale
\begin{equation}\tag{11.0.1}
\Phi_o=A_o/A_o^{\,o}
\end{equation}
ya considerado en la sección 1. Para todo número primo $\ell$,
denotemos por
\begin{equation*}
\Phi_o(\ell)
\end{equation*}
su componente $\ell$-primaria. Así, conocer $\Phi_o$ equivale a conocer
los $\Phi_o(\ell)$, en virtud de la fórmula
\begin{equation}\tag{11.0.2}
\Phi_o=\coprod_\ell\Phi_o(\ell).
\end{equation}

%%%%%%%%%%  scan idx 460  |  folio 449  %%%%%%%%%%

Obsérvese que los grupos $\Phi_o$ y $\Phi_o(\ell)$ están determinados,
respectivamente, por los grupos finitos ordinarios $\Phi_o(\bar k)$ y
$\Phi_o(\ell)(\bar k)$, junto con la acción natural de
\begin{equation*}
\pi_o=\operatorname{Gal}(\bar k/k)
\end{equation*}
sobre ellos. Por otra parte, está claro que $\Phi_o(\bar k)$ es también
el grupo de componentes conexas del modelo de Néron de
$A_{\widetilde K}$, donde $\widetilde K$ es el cuerpo de fracciones de
la henselización estricta de $S$ relativa a la clausura separable
elegida $\bar k$ de $k$. Por tanto, para estudiar $\Phi_o$, podemos
reducirnos al caso en que $S$ es estrictamente local; las acciones de
\begin{equation*}
\pi_o\simeq\operatorname{Gal}(\widetilde K/K)
\end{equation*}
sobre $\Phi_o(\bar k)$ se recuperan entonces automáticamente por
transporte de estructura. Del mismo modo, podríamos reducirnos además
al caso en que $S$ es completo, mediante el cambio de base
$\widehat S\longrightarrow S$ y el hecho de que la formación del modelo
de Néron conmuta con él [21, teorema 3.4 b)].

\smallskip
\noindent
11.1. Supongamos, pues, que $S$ es estrictamente local, de modo que
$\Phi_o$ puede identificarse con el grupo finito ordinario
$\Phi_o(k)$. Se obtiene
\begin{equation}\tag{11.1.1}
\Phi_o(k)\simeq A_o(k)/A_o^o(k)
 \simeq A(S)/A^o(S)
 \simeq A(K)/A(K)^o,
\end{equation}
donde la primera igualdad se deduce de que $A_o$ es liso sobre $k$; la
segunda, del lema de Hensel, pues $A$ es liso sobre $S$ y $S$ es
henseliano; y la tercera, de la relación
$A(S)\xrightarrow{\sim}A(K)$ y de la convención
\begin{equation}\tag{11.1.2}
A(K)^o\mathrel{\mathop:}=A^o(S).
\end{equation}
Así, $A(K)^o$ aparece como un subgrupo natural de índice finito en
$A(K)$ (y no debe confundirse con $A^o(K)=A(K)$). Los isomorfismos

%%%%%%%%%%  scan idx 461  |  folio 450  %%%%%%%%%%

(11.1.1) son evidentemente funtoriales en $A$ y compatibles con el
transporte de estructura en $(S,A_K)$ y, en particular, con la acción
de cualesquiera grupos de Galois.

Sea ahora $\ell$ un número primo distinto de $p$. Entonces $A^o$ es
$\ell$-divisible; más precisamente, $\ell\operatorname{id}_{A^o}$ es
un morfismo étale y sobreyectivo. En consecuencia, induce un
endomorfismo sobreyectivo de $A^o(S)=A(K)^o$, que es por tanto un
\underline{grupo $\ell$-divisible}. Se sigue inmediatamente que, para
todo entero $n\geq0$ primo con $p=\operatorname{char}(k)$,
\begin{equation}\tag{11.1.3}
{}_n\Phi_o(k)\simeq{}_nA(K)/{}_nA(K)^o,
 \qquad (n,p)=1.
\end{equation}
Por la definición (2.2.3) de la parte fija de ${}_nA_K$, se tiene
\begin{equation}\tag{11.1.4}
{}_nA(K)^o=({}_nA_K)^f(\bar K)
 \simeq T_n(A_K)^f(\bar K)
 \otimes_{\mathbb Z_n}\mathbb Z/n\mathbb Z.
\end{equation}
Tomando para $n$ una potencia de $\ell$ y poniendo, como en (2.5.3),
\begin{equation}\tag{11.1.5}
U=T_\ell\bigl(A(\bar K)\bigr),
 \qquad
T_\ell\bigl(A(\bar K)\bigr)^f=U^I,
\end{equation}
donde $U^I$ denota el submódulo de $U$ formado por los invariantes bajo
el grupo de inercia $I$, podemos escribir (11.1.3) en la forma
\begin{equation}\tag{11.1.6}
{}_n\Phi_o(k)\simeq (U_n)^I/(U^I)_n,
 \qquad n=\ell^{r+1},
\end{equation}
donde el subíndice $n$ denota la reducción módulo $n$, es decir, el
producto tensorial sobre $\mathbb Z_\ell$ con
$\mathbb Z_\ell/n\mathbb Z_\ell$. Como los morfismos de transición
entre los grupos que aparecen en ambos miembros de (11.1.6) son
evidentes, al pasar al límite inductivo sobre las potencias de $\ell$
se obtiene el

%%%%%%%%%%  scan idx 462  |  folio 451  %%%%%%%%%%

resultado siguiente.

\medskip
\noindent
\underline{Proposición 11.2.} \underline{Supongamos que $S$ es
estrictamente local y sea $\ell$ un número primo distinto de $p$.
Entonces existe un isomorfismo canónico}
\begin{equation}\tag{11.2.1}
\Phi_o(\ell)\simeq
 \frac{(U\otimes D_\ell)^I}{U^I\otimes D_\ell},
\end{equation}
\underline{donde $U$ se define en (11.1.5) y
$D_\ell=\mathbb Q_\ell/\mathbb Z_\ell$.}

Así, el grupo $\Phi_o(\ell)$ puede describirse únicamente en términos
del módulo de Tate $U=T_\ell(A(\bar K))$, considerado como módulo de
Galois bajo la acción del grupo de inercia $I$.
